\subsection{Kiểm thử}
Phần này trình bày kết quả kiểm thử của hệ thống, bao gồm các test case đã được thực hiện, kết quả đạt được và các lỗi (nếu có) đã được phát hiện trong quá trình kiểm thử. Các test case được thiết kế dựa trên các use case đã xác định trong giai đoạn phân tích hệ thống, nhằm đảm bảo rằng tất cả các chức năng của hệ thống đều hoạt động đúng như mong đợi. Kết quả kiểm thử sẽ được tổng hợp và đánh giá để xác định mức độ hoàn thiện và chất lượng của hệ thống trước khi triển khai chính thức.

Kiểm thử Module 1: Lập kế hoạch du lịch đã được thực hiện với các test case chi tiết được mô tả trong phần tiếp theo. Các test case này bao gồm cả các trường hợp thành công và các trường hợp lỗi để đảm bảo rằng hệ thống có thể xử lý tốt các tình huống khác nhau mà người dùng có thể gặp phải khi sử dụng chức năng lập kế hoạch du lịch.

\hypertarget{target:tc_m1}{}
\begin{landscape}
\small 
\begin{longtable}{|p{2.2cm}|p{2.2cm}|p{4cm}|p{4.5cm}|p{6cm}|p{2cm}|}
    
    % --- PHẦN 1: TIÊU ĐỀ CHO TRANG ĐẦU TIÊN ---
\caption{Danh sách Test Case chi tiết - Module 1: Lập kế hoạch du lịch} \label{tab:tc_module1} \\
\hline
\textbf{ID} & \textbf{ Use case} & \textbf{Tên Test case} & \textbf{Dữ liệu kiểm thử} & \textbf{Kết quả mong đợi} & \textbf{Kết quả } \\
\hline
\endfirsthead

% --- PHẦN 2: TIÊU ĐỀ CHO CÁC TRANG TIẾP THEO ---
\multicolumn{6}{c}{{\bfseries \tablename\ thetable{}: Danh sách Test Case chi tiết - Module 1: Lập kế hoạch du lịch}} \\
\hline
\textbf{ID} & \textbf{ Use case} & \textbf{Tên Test case} & \textbf{Dữ liệu kiểm thử} & \textbf{Kết quả mong đợi} & \textbf{Kết quả } \\
\hline
\endhead

\hline
\endlastfoot

% ================= NỘI DUNG BẢNG =================

% --- LUỒNG CHÍNH ---
TC-M1-01 & UC-M1-01 & Tạo lộ trình thành công & 
- Điểm đến: Đà Lạt \newline 
- Ngày: 01/12 - 03/12 \newline 
- Ngân sách: 4.000.000 \newline 
- Đi solotrip \newline
- Sở thích: Foodtour & 
Hệ thống hiển thị bản lộ trình nháp (Timeline view) với các địa điểm phù hợp tiêu chí đã chọn. & Đạt\\
\hline

TC-M1-02 & UC-M1-01 & Kiểm tra lỗi Ngày kết thúc nhỏ hơn Ngày bắt đầu & 
- Start date: 03/12 \newline
- End date: 01/12 & 
Hệ thống không cho phép chọn end date nhỏ hơn start date và yêu cầu nhập lại. & Đạt \\
\hline

TC-M1-03 & UC-M1-01 & Ràng buộc không hợp lý & 
- Ngân sách: 50.000 \newline
- Ngày: 01/12 - 03/12 & 
Hệ thống thông báo: "Failed to create matching itinerary" & Đạt\\
\hline

TC-M1-04 & UC-M1-02 & Lưu lộ trình thành công & 
N/A & 
Hệ thống thông báo "Save successful." & Đạt \\
\hline

TC-M1-05 & UC-M1-02 & Lưu lộ trình trùng tên& 
N/A & 
Hệ thống thông báo lỗi: "Existing itinerary with same name found" & Đạt \\
\hline

TC-M1-06 & UC-M1-03 & Xoá lộ trình thành công & 
N/A & 
Hệ thống yêu cầu xác nhận xoá, sau khi xác nhận, lộ trình bị xoá khỏi danh sách. & Đạt \\
\hline

TC-M1-07 & UC-M1-04 & Huỷ thao tác xoá lộ trình & 
N/A & 
Hệ thống không thực hiện thao tác xoá và quay lại trạng thái trước đó. & Đạt \\
\hline


\end{longtable}
\end{landscape}

Chi tiết kiểm thử cho các module khác được trình bày trong phần phụ lục. Mỗi test case được mô tả chi tiết về các bước thực hiện, dữ liệu đầu vào và kết quả mong đợi, cùng với kết quả thực tế sau khi kiểm thử để đánh giá hiệu quả của hệ thống.
\begin{itemize}
    \item \hyperlink{target:tc_m2}{Danh sách Test Case chi tiết cho Module 2}
    \item \hyperlink{target:tc_m3}{Danh sách Test Case chi tiết cho Module 3}
    \item \hyperlink{target:tc_m4}{Danh sách Test Case chi tiết cho Module 4}
    \item \hyperlink{target:tc_m5}{Danh sách Test Case chi tiết cho Module 5}
    \item   \hyperlink{target:tc_m8}{Danh sách Test Case chi tiết cho Module 8}
\end{itemize}
